\chapter{Farmer Health}

\section*{Definition and Overview}
Farmer health concerns the particular physical and mental health risks faced by people who work in agriculture, and the measures that protect them. Farming is among the most hazardous occupations in Australia: it combines heavy machinery, large animals, chemicals, dust, solar exposure, and isolation, all in a physically demanding environment. Farmers also carry a disproportionate burden of mental illness and suicide compared with the general population. Because farms are often workplaces, homes, and communities at once, health risks spill across every aspect of a farming family's life. This chapter outlines the main health hazards of agriculture and the practical strategies, services, and support networks that keep farming people safe and well.

\section*{Epidemiology}
Agriculture employs around 2--3\% of the Australian workforce yet accounts for a far higher share of workplace fatalities and serious injuries. Farming consistently ranks among the occupations with the highest rates of work-related death, with tractor and quad-bike incidents, being crushed or struck by animals, and drowning among the leading causes. Farmers have elevated rates of hearing loss, skin cancer, musculoskeletal injury, and respiratory disease, and report high levels of stress. Perhaps most concerning, suicide rates among male farmers are significantly above the national average, and financial pressures, drought, and isolation are recognised risk factors. Mental health support services, including dedicated rural and agricultural programs, have been established in response.

\section*{Aetiology and Pathophysiology}
The health risks of farming arise from the interaction of the environment, equipment, and demands of the work.
\begin{itemize}
    \item \textbf{Machinery and vehicles:} tractors, quad bikes, and other equipment cause rollovers, crush injuries, entrapment, and amputation, often involving children and older people on the property
    \item \textbf{Animals:} cattle, horses, and other livestock can crush, kick, gore, or trample handlers, causing blunt and penetrating trauma
    \item \textbf{Chemicals:} pesticides, herbicides, fuels, and veterinary products can be inhaled, absorbed through skin, or ingested, causing acute poisoning or long-term health effects
    \item \textbf{Dusts and moulds:} grain dust, hay dust, and moulds cause asthma, hypersensitivity pneumonitis (farmer's lung), and chronic airway disease
    \item \textbf{Noise:} prolonged exposure to engines, augers, and other machinery causes irreversible noise-induced hearing loss
    \item \textbf{Solar exposure:} outdoor work produces very high cumulative ultraviolet exposure and skin cancer risk
    \item \textbf{Psychosocial stress:} financial uncertainty, weather and commodity dependence, long hours, and social and geographic isolation drive anxiety, depression, and suicide risk
    \item \textbf{Access barriers:} distance from services, long working hours, and reluctance to seek help delay care and worsen outcomes
\end{itemize}

\section*{Clinical Features}
Farming-related health problems present in many forms, depending on the hazard.
\begin{itemize}
    \item Traumatic injuries: fractures, crush injuries, amputations, head injuries, and lacerations from machinery, animals, and falls
    \item Hearing loss: difficulty hearing conversation, tinnitus, and a history of noisy work without protection
    \item Skin problems: actinic keratoses, basal and squamous cell carcinomas, and melanoma on sun-exposed skin
    \item Respiratory symptoms: cough, wheeze, breathlessness, and flu-like illness after exposure to mouldy hay or grain dust
    \item Chemical exposure symptoms: headache, nausea, dizziness, skin irritation, and, in severe poisoning, seizures or collapse
    \item Mental health difficulties: persistent low mood, anxiety, irritability, sleep disturbance, harmful alcohol use, and suicidal thoughts
    \item Musculoskeletal strain: chronic back, shoulder, and knee pain from lifting, repetitive work, and prolonged machinery use
\end{itemize}

\section*{Differential Diagnosis}
Because farmers are exposed to many hazards at once, symptoms need careful sorting.
\begin{itemize}
    \item Chest symptoms could reflect farmer's lung, occupational asthma, infections, or cardiac disease
    \item Fatigue and headache may be from stress, dehydration, chemical exposure, carbon monoxide, or common illnesses
    \item Skin lesions require assessment to distinguish benign sun damage from skin cancer
    \item Low mood and anxiety may be primary mental illness or part of adjustment to financial and seasonal stress
    \item Injury patterns from farm machinery should always be assessed for head and internal injury
    \item Hearing problems may be occupational noise damage or unrelated middle and inner ear disease
\end{itemize}

\section*{When to Seek Medical Care}
Farmers and their families should not delay seeking care for injuries, persistent symptoms, or mental health concerns, and should plan ahead because of distance from services. Anyone experiencing suicidal thoughts should talk to a doctor, counsellor, or a helpline immediately.

\textbf{Call 000 (or local emergency services) for:}
\begin{itemize}
    \item Any serious machinery, animal, or vehicle injury, especially with severe bleeding, head injury, or suspected spinal injury
    \item Collapse, difficulty breathing, or fitting, which may indicate chemical poisoning or carbon monoxide exposure
    \item Chest pain or sudden breathlessness
    \item Any person found unconscious in or near machinery, silos, or confined spaces
\end{itemize}

\section*{Diagnosis and Investigations}
Diagnosis depends on the presenting problem but commonly involves:
\begin{itemize}
    \item \textbf{History:} detailed exposure history including machinery, animals, chemicals, dusts, and recent stressful events
    \item \textbf{Physical examination:} assessment of injuries, hearing, skin, lungs, and mental state
    \item \textbf{Lung function tests:} spirometry for respiratory symptoms, with specific testing for occupational asthma or farmer's lung
    \item \textbf{Audiometry:} hearing tests to quantify noise damage and guide protection
    \item \textbf{Blood and urine tests:} for suspected chemical exposure, infection, or organ injury
    \item \textbf{Skin examination:} full skin checks by a doctor familiar with sun damage
    \item \textbf{Mental health assessment:} structured questionnaires and discussion of stress, mood, alcohol, and suicide risk
\end{itemize}

\section*{Management}
\begin{itemize}
    \item \textbf{First aid and emergency care:} prompt treatment of injuries and rapid retrieval to hospital when needed
    \item \textbf{Injury rehabilitation:} physiotherapy, occupational therapy, and graduated return to farm work after serious injury
    \item \textbf{Hearing protection and rehabilitation:} earplugs or muffs for noisy work, with hearing aids where damage is established
    \item \textbf{Skin cancer care:} regular skin checks and removal of suspicious lesions, with sunscreen and protective clothing as the mainstay
    \item \textbf{Respiratory protection:} masks and dust control for grain and hay handling, and treatment for asthma and farmer's lung
    \item \textbf{Mental health care:} counselling, psychological therapy, medication where appropriate, and practical financial and business support
    \item \textbf{Preventive health:} routine check-ups, immunisations including tetanus, and management of blood pressure, cholesterol, and diabetes
\end{itemize}

\section*{Complications}
\begin{itemize}
    \item Disabling and fatal injuries from machinery rollovers, quad-bike accidents, and animal attacks
    \item Chronic, irreversible hearing loss and tinnitus
    \item Skin cancers, including melanoma, which can be fatal if detected late
    \item Progressive lung disease from repeated dust and mould exposure
    \item Chronic pain and disability from musculoskeletal injury, limiting ability to work
    \item Depression, harmful alcohol use, and suicide in the setting of unrelieved stress and isolation
    \item Financial hardship and loss of the farm when a breadwinner is injured or unwell
\end{itemize}

\section*{Prognosis and Outcomes}
Most farming injuries are preventable, and the prognosis for treated injuries is good when care is reached promptly. Chronic conditions such as hearing loss and skin cancer can be largely prevented by consistent protection, and lung disease risk falls when exposure is controlled. Mental health problems respond well to treatment, especially when farmers seek help early, and dedicated rural support services have saved many lives. The strongest predictor of good outcomes across all farming health risks is a culture that prioritises safety, protects against known hazards, and treats mental health as seriously as physical health.

\section*{Prevention and Self-Care}
\begin{itemize}
    \item Fit rollover protection to tractors, ride quad bikes only when trained, and never carry passengers
    \item Keep children away from machinery, chemicals, dams, and livestock handling areas
    \item Wear hearing protection, dust masks, sunscreen, and appropriate clothing for the task
    \item Read chemical labels, use personal protective equipment, and store chemicals securely out of reach
    \item Maintain a first-aid kit, know emergency numbers, and have a plan for getting help in remote areas
    \item Take regular breaks, protect sleep, and share workload with family or workers during busy seasons
    \item Talk openly about stress and finances, and use services such as farm counsellors, the Rural Adversity Mental Health Program, and helplines
    \item Have regular medical check-ups, skin checks, and hearing tests
\end{itemize}

\section*{Sources and Further Reading}
The following reputable sources provide further information for healthcare students and patients seeking reliable, current advice about farmer health.
\begin{itemize}
    \item National Centre for Farmer Health. \textit{Farmer health resources.} https://www.farmerhealth.org.au/
    \item Healthdirect Australia. \textit{Rural health services and farm safety.} https://www.healthdirect.gov.au/
    \item Safe Work Australia. \textit{Agriculture industry safety.} https://www.safeworkaustralia.gov.au/
    \item beyondblue and RU OK. \textit{Mental health support for farmers.} https://www.beyondblue.org.au/
    \item Mayo Clinic. \textit{Occupational health and injury prevention.} https://www.mayoclinic.org/
    \item World Health Organization. \textit{Occupational health and safety.} https://www.who.int/
\end{itemize}
